\chapter{Le second degré dans toutes ses formes}
\section*{Exercice 3.1 -- Choisir une forme}
Pour $f(x)=2x^2-8x+6$, donner les formes canonique et factorisée, puis son minimum.

\section*{Exercice 3.2 -- Trois discriminants}
Résoudre $x^2-6x+9=0$, $2x^2+x-3=0$ et $x^2+x+1=0$ dans $\R$.

\section*{Exercice 3.3 -- Inéquation factorisée}
Résoudre $-2(x+1)(x-4)\ge0$.

\section*{Exercice 3.4 -- Deux racines imposées}
Trouver tous les polynômes de degré exactement $2$ s'annulant en $-2$ et $3$. Lequel vaut $12$ en $0$ ?

\section*{Exercice 3.5 -- Paramètre}
Selon le réel $m$, combien de solutions possède $x^2-2x+m=0$ ?

\section*{Exercice 3.6 -- Optimisation géométrique}
Un rectangle a un périmètre de $28$ cm. Montrer que son aire est au plus $49$ cm$^2$ et préciser l'égalité.
